\part{Cómo funciona LoRa}
\label{part:lora}

\chapter{Qué es modular una señal}

\section{Qué es}

Una antena por sí sola solo puede emitir una onda a una frecuencia fija; para transportar
información hay que \textbf{modular}: variar alguna propiedad de esa onda (amplitud,
frecuencia o fase) de forma que el receptor pueda reconstruir los bits originales a partir de
esas variaciones. AM varía la amplitud, FM varía la frecuencia; los esquemas digitales modernos
—incluido LoRa— varían combinaciones de frecuencia y fase para codificar símbolos, no solo bits
sueltos.

\section{Por qué importa aquí}

LoRa no usa un esquema de modulación clásico como FSK o PSK puros: usa \textbf{espectro
ensanchado por chirp} (\emph{Chirp Spread Spectrum}, CSS), diseñado específicamente para
maximizar el alcance a costa de velocidad, que es exactamente el compromiso que necesita un
sensor remoto que envía pocos bytes por segundo pero debe llegar lejos con un transmisor de
baja potencia.

\chapter{Espectro ensanchado por chirp}

\section{Qué es}

Un \emph{chirp} es un tono cuya frecuencia sube (o baja) linealmente en el tiempo, barriendo
todo el ancho de banda del canal durante la duración de un símbolo. Es el mismo fenómeno que un
silbido que empieza grave y termina agudo.

\begin{nota}[title={Analogía}]
Imagina que en vez de gritar una sola nota fija (como haría una modulación de amplitud
clásica), silbas deslizando el tono de grave a agudo. Ese deslizamiento completo, con una
velocidad de barrido determinada, \emph{es} un símbolo. LoRa codifica información desplazando
el \emph{punto de partida} de ese barrido: un silbido que empieza a medio camino y da la vuelta
codifica un símbolo distinto de uno que empieza al principio.
\end{nota}

\begin{figure}[h]
  \centering
  \begin{tikzpicture}[scale=1.0, font=\footnotesize]
    \draw[-{Latex}] (0,0) -- (6.5,0) node[right] {tiempo};
    \draw[-{Latex}] (0,0) -- (0,3.2) node[above] {frecuencia};
    % up-chirp completo (símbolo 0)
    \draw[thick, blue!60!black] (0,0.3) -- (2,3);
    \draw[thick, blue!60!black, dashed] (2,3) -- (2,0);
    \draw[thick, blue!60!black] (2,0) -- (4,2.7);
    \node[blue!60!black, below, text width=2.6cm, align=center] at (1.6,-0.15)
      {símbolo 0\\(arranca en 0)};
    % chirp con offset (símbolo desplazado)
    \draw[thick, red!60!black] (4,1.6) -- (4.9,3);
    \draw[thick, red!60!black, dashed] (4.9,3) -- (4.9,0);
    \draw[thick, red!60!black] (4.9,0) -- (6,1.6);
    \node[red!60!black, below, text width=2.4cm, align=center] at (5.3,-0.15)
      {símbolo\\con offset};
  \end{tikzpicture}
  \caption{Un chirp barre linealmente toda la banda; el punto donde empieza el barrido
    (y da la vuelta) codifica el valor del símbolo. Cada línea discontinua marca un
    ``salto'' de vuelta a la frecuencia mínima al llegar al máximo del canal.}
\end{figure}

Por qué esto ayuda al alcance: el receptor no necesita detectar un chirp aislado con precisión
absoluta de amplitud (como sí exige FSK/PSK clásicos); necesita \emph{correlacionar} la señal
recibida con la forma conocida del chirp. Esa correlación acumula energía de todo el símbolo, lo
que permite decodificar señales que están \emph{por debajo del nivel de ruido} —algo imposible
con modulaciones de banda estrecha convencionales—.

\section{Factor de dispersión: el compromiso alcance/velocidad}

\subsection*{Qué es}

El \textbf{factor de dispersión} (\emph{Spreading Factor}, SF) determina cuántos \emph{bits por
símbolo} se transmiten y, por tanto, cuánto dura cada chirp: $2^{SF}$ posibles puntos de
arranque distintos, es decir, $SF$ bits de información por símbolo. Un SF más alto significa
más puntos de arranque posibles (más información por símbolo) pero también un barrido más
lento, así que cada símbolo tarda más tiempo en transmitirse.

\begin{equation}
  T_{\text{sym}} = \frac{2^{SF}}{BW}
  \label{eq:tsym}
\end{equation}

\subsection*{Por qué importa aquí}

Un SF más alto extiende la energía de cada bit sobre un chirp más largo, lo que mejora la
\textbf{sensibilidad} del receptor (puede detectar señales más débiles) a costa de una
\textbf{tasa de datos} menor. Este proyecto usa \textbf{SF7}, el más bajo (y más rápido) del
rango típico (7 a 12), porque el presupuesto de enlace de la sección~\ref{sec:presupuesto}
mostró un margen de \SI{45}{dB} incluso con SF7: no hace falta sacrificar velocidad para ganar
alcance que ya sobra.

\begin{table}[h]
  \centering
  \small
  \begin{tabular}{@{}lrrl@{}}
    \toprule
    SF & Símbolos/s @ 500 kHz & Sensibilidad aprox. & Compromiso \\
    \midrule
    7  & \num{3906}  & $\approx \SI{-118}{dBm}$ & rápido, alcance moderado \\
    9  & \num{977}   & $\approx \SI{-124}{dBm}$ & más lento, más alcance \\
    12 & \num{122}   & $\approx \SI{-131}{dBm}$ & muy lento, máximo alcance \\
    \bottomrule
  \end{tabular}
  \caption{A mayor SF, mayor sensibilidad (más alcance) pero menor tasa de símbolos
    (más tiempo de aire por trama). Valores de sensibilidad aproximados para BW = 500 kHz.}
\end{table}

\section{Ancho de banda y corrección de errores}

\subsection*{Qué es}

El \textbf{ancho de banda} (BW) es el rango de frecuencias que ocupa cada chirp: cuanto más
ancho, más rápido se puede barrer (menos tiempo por símbolo, ecuación~\eqref{eq:tsym}), pero
también entra más ruido de banda ancha al receptor, lo que reduce ligeramente la sensibilidad.
Este proyecto usa \textbf{\SI{500}{kHz}}, el máximo disponible en el módulo, para minimizar el
tiempo de aire.

La \textbf{tasa de codificación} (\emph{Coding Rate}, CR) añade bits redundantes de corrección
de errores hacia adelante (FEC): con CR = 4/5, por cada 4 bits de datos se transmite 1 bit de
paridad. Esto permite al receptor corregir errores de bit sin pedir retransmisión, a costa de
transmitir algo más de información por trama útil. El proyecto usa \textbf{4/5}, la tasa mínima
de sobrecarga (frente a 4/6, 4/7 o 4/8, más robustas pero más lentas).

\subsection*{Por qué importa aquí}

BW alto y CR mínimo son ambos decisiones de \emph{minimizar el tiempo de aire}, coherentes con
la elección de SF7: el margen de enlace de \SI{45}{dB} permite ser agresivo en velocidad sin
poner en riesgo la fiabilidad del enlace a \SI{600}{m}.

\chapter{El tiempo de aire, término a término}
\label{sec:airtime}

\section{Qué es}

El \textbf{tiempo de aire} (\emph{time on air}) es cuánto dura, en el aire, la transmisión
completa de una trama: preámbulo más carga útil. Es el número que determina cuántas tramas por
segundo caben en el canal y, por extensión, cuántos nodos puede soportar el sistema —el tema de
la Parte~\ref{part:regulacion} y del diseño TDMA de la Fase~4—.

La fórmula estándar de LoRa (documentada por Semtech) tiene dos partes: el preámbulo y la carga
útil, cada una expresada en número de símbolos y luego convertida a tiempo con $T_{\text{sym}}$
de la ecuación~\eqref{eq:tsym}.

\subsection*{Preámbulo}

\begin{equation}
  T_{\text{preámbulo}} = (n_{\text{pre}} + 4.25) \cdot T_{\text{sym}}
  \label{eq:tpre}
\end{equation}

$n_{\text{pre}}$ es el número de símbolos de preámbulo configurados (típicamente 8: una
secuencia conocida que el receptor usa para sincronizarse). El término fijo $4.25$ cubre la
palabra de sincronización y el arranque de la decodificación, y es parte de la especificación
de Semtech, no un parámetro configurable.

\subsection*{Carga útil}

\begin{equation}
  n_{\text{payload}} = 8 + \max\!\left(
    \left\lceil \frac{8 \cdot PL - 4 \cdot SF + 28 + 16}{4 \cdot SF} \right\rceil
    \cdot (CR + 4),\ 0 \right)
  \label{eq:npayload}
\end{equation}

\begin{equation}
  T_{\text{trama}} = T_{\text{preámbulo}} + n_{\text{payload}} \cdot T_{\text{sym}}
  \label{eq:ttrama}
\end{equation}

donde $PL$ es el tamaño de la carga útil en bytes, y $CR$ es la tasa de codificación expresada
como el número de bits de paridad por cada 4 bits de datos (CR = 1 para 4/5, CR = 4 para 4/8).
El término $+16$ dentro del numerador corresponde al CRC de 16 bits que valida la integridad de
la trama a nivel de radio (independiente del CRC16 propio de la capa de aplicación, que se
implementa en la Fase~2); el término $+28$ y el $8$ inicial son constantes de la
especificación LoRa relacionadas con el encabezado explícito.

\section{Por qué importa aquí: los 76,9 ms resueltos}

El proyecto agrupa 10 muestras por trama. Con el formato de trama de la Fase~2 (cabecera +
10 muestras de \SI{18}{B} + CRC16, ver Supuestos declarados en el plan), el tamaño resultante
es $PL = \SI{194}{B}$. Con SF7, BW = \SI{500}{kHz} y CR = 1 (4/5):

\begin{ejemplo}
\textbf{Paso 1 — Duración de un símbolo:}
\[
  T_{\text{sym}} = \frac{2^7}{500\,000} = \SI{0.256}{ms}
\]

\textbf{Paso 2 — Preámbulo (8 símbolos configurados):}
\[
  T_{\text{preámbulo}} = (8 + 4.25) \cdot 0.256\,\text{ms} = \SI{3.136}{ms}
\]

\textbf{Paso 3 — Símbolos de carga útil, con $PL = 194$, $SF = 7$, $CR = 1$:}
\begin{align*}
  \frac{8 \cdot 194 - 4 \cdot 7 + 28 + 16}{4 \cdot 7}
    &= \frac{1552 - 28 + 28 + 16}{28} = \frac{1568}{28} = 56 \\
  n_{\text{payload}} &= 8 + \lceil 56 \rceil \cdot (1 + 4) = 8 + 56 \cdot 5 = 288
\end{align*}

\textbf{Paso 4 — Tiempo de trama:}
\begin{align*}
  T_{\text{trama}} &= 3.136\,\text{ms} + 288 \cdot 0.256\,\text{ms} \\
    &= 3.136 + 73.728 = \SI{76.864}{ms} \approx \SI{76.9}{ms}
\end{align*}
\end{ejemplo}

Este es exactamente el valor que reproduce \texttt{common/airtime.py} en la Fase~1, y el que
sostiene el diseño del slot TDMA de \SI{136.9}{ms} (76,9 ms de trama + 60 ms de guarda) en la
Fase~4.

\section{Por qué un SF alto no siempre es mejor}

Es tentador pensar que, si SF7 ya funciona con \SI{45}{dB} de margen, usar SF12 —el más
robusto— sería ``más seguro''. En este proyecto sería contraproducente: la
ecuación~\eqref{eq:tsym} muestra que $T_{\text{sym}}$ crece \emph{exponencialmente} con SF, no
linealmente. Pasar de SF7 a SF12 multiplica $T_{\text{sym}}$ por $2^5 = 32$, lo que dispararía
el tiempo de aire de cada trama muy por encima de lo que 30 nodos pueden compartir en un canal
sin colisionar constantemente (ver el diseño de superframe de la Fase~4). El SF correcto no es
el más robusto en abstracto, sino el más bajo que todavía deja margen de enlace suficiente para
la distancia real del despliegue —exactamente el criterio que usó la sección~\ref{sec:presupuesto}
para justificar SF7—.

\section{Dónde aparece en el código}

\texttt{common/airtime.py} (Fase~1) implementa las ecuaciones~\eqref{eq:tsym},
\eqref{eq:tpre}, \eqref{eq:npayload} y \eqref{eq:ttrama} literalmente, y su prueba
(\texttt{tests/test\_airtime.py}) verifica el resultado de \SI{76.9}{ms} calculado aquí. El
mismo módulo calcula la duración de slot y la capacidad del canal que sostienen el diseño
TDMA de la Fase~4.
